\section{Scope and retrospective audit design}
\label{sec:design}

We audited claims from four campaigns in one split language model research
programme, then checked two broader claims about security scope. A claim
and its supporting records form the unit of analysis. We selected the
campaigns used in the original synthesis because they let us trace how
individual results became broader conclusions. The selection serves that
purpose, but provides no estimate of how often privacy evaluations fail
across the field. Each component result remains attributed to its source
campaign. The audit reconciles those records with the conclusions they
support.

The historical outcomes were available before we specified this audit.
We distinguish frozen exploration, completed experiments, proposed
confirmation and the present secondary analysis throughout. Our work
recomputes stored arithmetic and examines existing evidence without
retraining victim models, collecting captures, fitting stronger attackers
or filling gaps with inferred outcomes. Independent replication and
prospective confirmation would require their own studies. The
LLM-Attack Suite supplies relevant evaluation contracts and
controls~\citep{suite2026}. Here we use that context to examine claims
across campaigns, with no new validation of the suite itself.

\paragraph{Evidence required for a claim.}
For each claim, we record the protocol and cohort, the adversary's access
to secrets and remote state, and the channels it can observe. We also
identify the protected target and its alignment with observations, the
attack and training budgets, and the score, baseline and acceptance rule.
The qualification binding and the independent unit used to calculate
uncertainty complete the record. A qualification binding specifies the
target, channel, defence condition and implementation covered by positive
and null controls. A control passed elsewhere in the programme supplies
no qualification for a new binding. Likewise, recording a channel alone
does not evaluate attacks against it.

Matching generated records and currently admitted manifests take
precedence over narrative summaries. When a later decision supersedes an
earlier one, we retain that history together with the limitations attached
to the current status. This rule applies whichever decision is more
favourable. The source index links each retained quantity to its artifact
and distinguishes new arithmetic from results read through an existing
audit or manuscript. These links establish traceability within the supplied
package. Completeness of the original collection remains a separate
question.

\paragraph{Five outcomes for a specified claim.}
Different endpoints of the same configuration can have different
outcomes. We use the following terms consistently.
\begin{description}
  \item[Breach.] An admissible, qualified measurement exceeds the
    declared privacy rule. It establishes a breach for the specified
    observer, target, attack and finite protocol, without necessarily
    quantifying every disclosure.
  \item[Non-detection.] An admissible, qualified test stays within its
    rule for the tested attack and budget. This outcome supplies neither
    an upper bound on leakage nor a privacy guarantee for the whole system.
  \item[Undetermined.] Missing coverage, incomplete qualification or
    unavailable admissible observations prevents the required conclusion.
    The endpoint cannot contribute a privacy pass.
  \item[Invalid.] A demonstrated defect disconnects the measurement from
    the claimed estimand. For example, labels for the next token cannot
    establish resistance to reconstruction of the current token. The
    original score may still describe another clearly named quantity.
  \item[Out-of-scope.] The claim requires an observer, access to secrets,
    target or protocol different from the one evaluated. Even an accurate
    measurement in the original setting cannot establish the changed claim.
\end{description}

We record utility and remote gain separately and preserve each campaign's
aggregation rule. In the privacy, utility and gain exploration, a qualified
privacy breach or an independent utility failure can reject a configuration
whose other privacy endpoints remain undetermined. Incomplete qualification
blocks a conjunctive pass without erasing valid negative evidence. Failure
to meet an exploratory gain rule also leaves the population gain uncertain.

\paragraph{Transferring a conclusion.}
A conclusion can carry to another setting when the relevant protocol,
observer access, observations, target and measurement binding match, or
when an explicit argument justifies the change. A result for a wire
observer therefore needs further justification before it can support a
claim about a compromised host. Empirical non-detection provides no
information theoretic ceiling, and correlated rows add no independent
training seeds. The cases illustrate these limits on inference. They
provide neither an independent validation cohort nor evidence for a
universal privacy decision procedure.
